% Lecture 06: Конструкторы, деструкторы и копирование

\section{Лекция 6. Конструкторы, деструкторы и копирование}
\label{sec:lecture-06}

В блоке~5 \texttt{CourseBlock::inspect} сначала создавал почти пустой объект,
а затем по одному заполнял private members. Доступ извне закрыт, но внутри
class всё ещё существует промежуток, когда объект не представляет настоящий
блок курса. Конструктор позволяет связать начало lifetime с установлением
корректного начального состояния.

\begin{importantbox}[Идея]
Запись \texttt{T b = a;} обычно создаёт новый объект \texttt{b}, а не второе
имя для \texttt{a}. У нового объекта свой lifetime и своё состояние. Чтобы
такая value semantics была корректной, тип должен правильно определять
создание, копирование, присваивание и уничтожение.
\end{importantbox}

\subsection{Создать объект --- значит начать его lifetime}

Class задаёт собственный тип, а объект этого типа хранит значение. Для
обычного local object определение выделяет подходящее storage, выполняет
инициализацию, и после успешного завершения инициализации начинается lifetime.
При выходе из scope lifetime заканчивается. Стандарт не требует описывать это
как размещение в конкретном CPU stack frame.

\begin{cppcode}
class CourseBlock {
public:
    CourseBlock(int number, const std::string& slug)
        : number_{number}, slug_{slug} {}

private:
    int number_;
    std::string slug_;
};

CourseBlock block{6, "constructors_destructors_and_copy"};
\end{cppcode}

\texttt{CourseBlock} внутри class --- имя constructor. У него нет return type:
даже \texttt{void} писать нельзя. Parameters задают данные, необходимые для
начала корректного состояния. Compiler проверяет существование подходящего
constructor, типы arguments и доступ. Runtime инициализирует subobjects
\texttt{number\_} и \texttt{slug\_}, затем выполняет body.

Здесь различим три записи:

\begin{cppcode}
Widget first;          // default initialization
Widget second{};       // value initialization
Widget third{42};      // direct-list-initialization
\end{cppcode}

Для class type первые две формы вызывают доступный default constructor;
детальные различия zero/value initialization важны прежде всего для типов без
пользовательского constructor и будут уточняться по мере необходимости.
Форма с argument напрямую выбирает подходящий constructor. Braces также
запрещают ряд narrowing conversions.

\subsubsection{Default constructor и слово explicit}

Constructor, который можно вызвать без arguments, называется default
constructor:

\begin{cppcode}
class Counter {
public:
    Counter() : value_{0} {}
private:
    int value_;
};
\end{cppcode}

Если class не объявляет constructors, compiler может неявно объявить default
constructor. Но как только programmer объявил constructor с параметрами,
автоматически получить ещё и parameterless constructor нельзя:

\begin{cppcode}
class BlockNumber {
public:
    explicit BlockNumber(int value) : value_{value} {}
private:
    int value_;
};

BlockNumber good{6};
// BlockNumber absent; // error: default constructor отсутствует
\end{cppcode}

\texttt{explicit} у one-argument constructor запрещает использовать его как
неявное преобразование \texttt{int -> BlockNumber}. Это естественно, когда
число не должно незаметно становиться domain object. Подробная overload
resolution здесь не нужна. Если пустое состояние осмысленно, default
constructor можно объявить явно или написать \texttt{Type() = default;}.
Добавлять его только ради удобства, разрушая invariant, не следует.

\subsection{Member initializer list: сразу создать parts}

После parameter list constructor стоит двоеточие и список mem-initializers:

\begin{cppcode}
CourseBlock::CourseBlock(int number, const std::string& slug)
    : number_{number}, slug_{slug} {
    // Здесь оба member object уже инициализированы.
}
\end{cppcode}

Это initialization, не assignment. Members создаются до входа в body. Вариант
«сначала создать пустую string, потом присвоить» выполнял бы другую
последовательность операций и вообще невозможен для reference member или
обычного \texttt{const} member:

\begin{cppcode}
class Selection {
public:
    Selection(const int& source, int limit)
        : source_{source}, limit_{limit} {}
private:
    const int& source_;
    const int limit_;
};
\end{cppcode}

Reference должна быть привязана, а const object --- получить значение при
инициализации. При этом reference не владеет referent: design обязан
гарантировать, что тот живёт достаточно долго.

\begin{warningbox}[Порядок задаёт class, а не список]
Non-static data members инициализируются в порядке их declarations в class.
Порядок записей после двоеточия этого не меняет. Поэтому mem-initializer list
пишут в том же порядке: иначе visual order обманывает reader, а зависимость
одного member от ещё не инициализированного может привести к ошибке.
\end{warningbox}

В \texttt{initializer\_order\_warning.cpp} class объявляет \texttt{first\_}
до \texttt{second\_}, но визуально перечисляет наоборот. GCC/Clang с
\texttt{-Wreorder} реально предупреждает, что фактический порядок иной.
CTest требует это warning; языковое правило действует независимо от warning.

При уничтожении содержащего объекта members уничтожаются в обратном порядке
завершения их construction. Так часть, созданная позже, исчезает раньше.

\subsection{Constructor как граница invariant}

Плохая модель заставляет caller выполнить протокол:

\begin{cppcode}
CourseBlock block;
block.set_number(6);
block.set_slug("constructors_destructors_and_copy");
block.rebuild_paths();
\end{cppcode}

Между строками существуют состояния «номер без slug» и «новая identity со
старыми paths». Private fields плюс public setters не устраняют проблему.
Лучше получить все обязательные данные в constructor и сразу вычислить
dependent members. Тогда доступный public interface не создаёт
полуинициализированный \texttt{CourseBlock}.

\begin{cppcode}
CourseBlock(const std::filesystem::path& root, int number,
            const std::string& full_name)
    : number_{number},
      full_name_{full_name},
      slug_{full_name.substr(3)},
      lecture_path_{root / "lectures" / (full_name + ".tex")},
      examples_path_{root / "examples" / full_name} {}
\end{cppcode}

Compiler не знает domain rule «номер от 1 до 15»; его проверяет
implementation. Полная стратегия сообщения об ошибочном construction будет
связана с exceptions позже. Сейчас experiment передаёт только известную
таблицу корректных identity.

\subsection{Destructor и конец lifetime}

Destructor --- special member function с именем \texttt{\textasciitilde Type}:

\begin{cppcode}
class Trace {
public:
    explicit Trace(const std::string& name) : name_{name} {
        std::cout << "construct " << name_ << '\n';
    }
    ~Trace() {
        std::cout << "destroy " << name_ << '\n';
    }
private:
    std::string name_;
};
\end{cppcode}

У destructor нет return type и parameters. Для обычного local object он
вызывается при выходе из scope; затем заканчивается lifetime. Перед
освобождением storage уничтожаются members. Вложенный scope лаборатории
делает порядок наблюдаемым, но не доказывает конкретную реализацию stack.

Destructor нужен типу для завершения владения ресурсом. Его не следует
вызывать вручную как обычную \texttt{cleanup()} function: явный вызов не
превращает object обратно в безопасное «пустое» значение, а последующее
автоматическое уничтожение способно нарушить lifetime rules. Ручное управление
lifetime относится к advanced low-level code, не к практике этого блока.

\begin{standardbox}[Deterministic destruction]
Для local objects с automatic storage duration обычный выход из scope
определяет момент destruction языковыми правилами. Это важная основа будущей
RAII, хотя полную теорию acquisition, exceptions и smart pointers мы изучим
позже.
\end{standardbox}

\subsection{Какие special members способен дать compiler}

В фокусе четыре операции:
default constructor; copy constructor; copy assignment operator; destructor.

При выполнении условий языка compiler неявно объявляет/определяет некоторые
из них. «Generated» не значит «ничего не делает»: операция вызывает
соответствующие операции members. Взаимодействие правил зависит от того, что
programmer объявил. Поэтому опасно писать один manual resource operation, не
проверив остальные. Move constructor и move assignment существуют, но их
семантика и Rule of Five относятся к будущему блоку.

\subsection{Copy construction: появляется новый объект}

Обычная форма copy constructor:

\begin{cppcode}
class Card {
public:
    Card(const Card& other)
        : number_{other.number_}, title_{other.title_} {}
private:
    int number_;
    std::string title_;
};
\end{cppcode}

Reference нужна, чтобы читать source, не пытаясь сначала скопировать его для
parameter. \texttt{const} разрешает копировать const source и обещает не
менять его через \texttt{other}. При

\begin{cppcode}
Card a{6, "copy"};
Card b = a;
\end{cppcode}

\texttt{b} ещё не существовал: это initialization, выбирающая copy
constructor, несмотря на знак \texttt{=}. У \texttt{b} отдельный lifetime и
его members копируются по copy semantics собственных типов. Copy
\texttt{std::string} и \texttt{std::vector} создаёт независимое value:
изменение элементов одной копии не меняет другую.

Pass-by-value создаёт parameter object из argument и поэтому может вызвать
copy. Return-by-value тоже может предполагать такую операцию на уровне
исходного кода, но фактический отдельный вызов не всегда происходит.

\subsubsection{Copy elision без подсчёта магических вызовов}

Copy elision позволяет не материализовать промежуточную копию. В C++17 есть
гарантированные случаи, где result object конструируется непосредственно в
назначении. В других случаях compiler может применить разрешённую elision.
Поэтому нельзя обещать: «каждый return class object вызывает один copy
constructor».

GCC/Clang flag \texttt{-fno-elide-constructors} полезен для отдельных
наблюдений, но это experimental toolchain flag, не другая семантика C++ и не
способ отменить гарантированные C++17 cases. Лаборатория проверяет
независимость values, а не хрупкое универсальное число сообщений.

\subsection{Copy assignment: объект уже существует}

\begin{cppcode}
class Card {
public:
    Card& operator=(const Card& other) {
        number_ = other.number_;
        title_ = other.title_;
        return *this;
    }
};
\end{cppcode}

Здесь compiler видит перегруженный assignment operator. Слева уже живой
object, его lifetime не начинается заново: состояние заменяется состоянием
\texttt{other}. Возврат \texttt{*this} по reference --- conventional
контрактом assignment и позволяет цепочки \texttt{a = b = c}; отдельный
объект для результата не создаётся.

\begin{cppcode}
Card b = a; // copy construction: b создаётся
b = a;      // copy assignment: b уже существует
\end{cppcode}

Self-assignment \texttt{b = b;} допустим. Generated assignment стандартных
members с ним справляется. Manual resource owner иногда использует
\texttt{if (this == \&other)}, чтобы не удалить resource до чтения из него.
Guard не обязателен для каждого корректного operator: можно
написать self-safe algorithm без отдельной ветви.

\subsection{Shallow copy ломает уникальное владение}

Рассмотрим намеренно старомодный учебный type:

\begin{cppcode}
class ShallowBuffer {
public:
    explicit ShallowBuffer(int value) : data_{new int{value}} {}
    ~ShallowBuffer() { delete data_; }
    int& value() { return *data_; }
private:
    int* data_; // owning raw pointer
};

ShallowBuffer first{10};
ShallowBuffer second = first; // generated copy копирует pointer value
\end{cppcode}

Здесь pointer --- не просто доступ, а заявление: object владеет allocation и
обязан освободить её. Generated copy выполняет memberwise copy и получает два
одинаковых адресных значения. Оба objects считают один resource своим.
Изменение через \texttt{second} видно через \texttt{first}; при destruction
первый \texttt{delete} освобождает allocation, второй пытается освободить его
ещё раз. Это double delete и undefined behavior, а не гарантированный crash.

\begin{definition}
\textbf{Shallow copy} resource-owning representation копирует handle/address,
не создавая независимый resource. \textbf{Deep copy} создаёт новый resource и
копирует его содержимое, сохраняя отдельное владение.
\end{definition}

Raw pointer также может быть non-owning view. Тогда он не освобождает referent,
но copy всё равно не продлевает lifetime. Ownership должен быть понятен из
design и документации. Naked \texttt{new}/\texttt{delete} в этом блоке
изолированы именно для демонстрации ошибки; domain model их не использует.

\subsection{Deep copy и Rule of Three}

Исправленный \texttt{DeepBuffer} определяет три операции:

\begin{cppcode}
DeepBuffer(const DeepBuffer& other)
    : size_{other.size_}, data_{new int[other.size_]} {
    for (std::size_t i = 0; i < size_; ++i) {
        data_[i] = other.data_[i];
    }
}

DeepBuffer& operator=(const DeepBuffer& other) {
    if (this == &other) {
        return *this;
    }
    int* replacement = new int[other.size_];
    for (std::size_t i = 0; i < other.size_; ++i) {
        replacement[i] = other.data_[i];
    }
    delete[] data_;
    data_ = replacement;
    size_ = other.size_;
    return *this;
}

~DeepBuffer() { delete[] data_; }
\end{cppcode}

Copy constructor выделяет отдельный array. Assignment сначала готовит замену,
потом освобождает старый resource; self-assignment guard не даёт читать уже
удалённый собственный array. Destructor освобождает ровно resource этого
object. Полный файл проверяет, что изменения копии независимы.

\begin{definition}[Rule of Three]
Если class вручную управляет resource и ему нужен custom destructor, copy
constructor или copy assignment, обычно необходимо внимательно рассмотреть
все три операции.
\end{definition}

Это design guideline, не синтаксическое требование стандарта. Compiler может
принять нарушение, как он принял \texttt{ShallowBuffer}; проблема проявится
runtime как неправильное ownership и UB.

\subsection{Rule of Zero: предпочтительный современный итог}

Ту же задачу хранения последовательности умеет решать member
\texttt{std::vector<int>}:

\begin{cppcode}
class Buffer {
public:
    Buffer(const std::string& label, const std::vector<int>& values)
        : label_{label}, values_{values} {}
private:
    std::string label_;
    std::vector<int> values_;
};
\end{cppcode}

Class не пишет destructor/copy constructor/copy assignment. String и vector
сами владеют ресурсами и корректно копируются; generated memberwise operations
дают независимое value. Это и есть практическая мотивация Rule of Zero:
строить обычный class из resource-managing members так, чтобы ему не требовался
ручной код special members.

Полная RAII и smart pointers будут позже. Уже сейчас вывод применим к
\texttt{CourseBlock}: members \texttt{int}, \texttt{bool},
\texttt{std::string}, \texttt{std::filesystem::path} и контейнеры имеют
подходящие операции. Добавлять custom copy/destructor «потому что лекция о
копировании» было бы лишним кодом и новой площадью для ошибок.

\subsection{Стоимость value semantics}

Generated не означает бесплатный:
copy scalar members обычно дёшев и имеет постоянную стоимость; copy string/vector может выделять storage и копировать characters или elements; стоимость class copy складывается из стоимости members; deep copy buffer длины $n$ требует $O(n)$ времени и отдельной памяти; destructor owner должен освободить resource; vector value semantics имеет цену, но даёт ясное независимое владение.

При copy assignment контейнер иногда переиспользует имеющуюся capacity.
Точные allocations зависят от library implementation и текущего состояния. Поэтому
не утверждаем конкретное число allocations без измерения. В этом блоке важна
семантика; BENCHMARKS --- N/A. Позже move semantics позволит в подходящих
сценариях передавать resource вместо deep copy, но её синтаксис пока не нужен.

\subsection{Эксперимент 1: copy\_semantics\_lab}

Каталог \texttt{01\_copy\_semantics\_lab} объединяет шесть связанных cases.
\texttt{LessonCard} проверяет copy construction и copy assignment обычного
value type. Копия vector меняет element и добавляет topic, не меняя source.
\texttt{Trace} печатает constructor, copy constructor, assignment и
destructor во вложенном scope. После строки \texttt{-- nested scope end --}
destructors идут в обратном порядке construction.

\texttt{DeepBuffer} проверяет Rule of Three, независимые allocations и
self-assignment. \texttt{ZeroBuffer} получает ту же value semantics от string
и vector без custom copy/destructor. Safe target должен закончить строкой
\texttt{COPY\_SEMANTICS\_LAB=PASS}.

\begin{shellcode}
cd examples/06_constructors_destructors_and_copy/01_copy_semantics_lab
cmake -S . -B build -DCMAKE_BUILD_TYPE=Debug
cmake --build build
./build/copy_semantics_lab
ctest --test-dir build --output-on-failure
\end{shellcode}

Bad shallow copy отделён opt-in option:

\begin{shellcode}
cmake -S . -B build-sanitize -DENABLE_UNSAFE_CASE=ON \
  -DCMAKE_BUILD_TYPE=Debug
cmake --build build-sanitize
./build-sanitize/bad_shallow_copy
\end{shellcode}

На GCC~13 в фактическом запуске process завершился с nonzero status, а ASan
сообщил \texttt{attempting double-free}. Это подтверждает выполненный bad
path на этой toolchain. Отсутствие sanitizer diagnostic в другом запуске не
доказывало бы отсутствия UB. Normal target намеренно не падает; его safe cases
чисты и под ASan/UBSan.

\subsection{Эксперимент 2: независимый inventory snapshot}

\texttt{02\_inventory\_snapshot} развивает class-based inventory блока~5,
не меняя его файлы. Constructor \texttt{CourseBlock} сразу устанавливает
number, identity, derived paths и измеренное состояние. Затем реальные данные
всех 15 блоков проходят две разные операции:

\begin{cppcode}
const std::vector<CourseBlock> inventory = collect_blocks(root, names);
std::vector<CourseBlock> inventory_copy = inventory; // construction
std::vector<CourseBlock> inventory_assigned;
inventory_assigned = inventory;                      // assignment
\end{cppcode}

Custom special members у \texttt{CourseBlock} отсутствуют. Сравнение
observable fields подтверждает, что generated copy сохраняет данные. Чтобы не
портить read-only domain interface искусственным setter, для анализа введён
отдельный \texttt{MutableBlockSnapshot}. Изменение label и vector notes в
копии не затрагивает original; assignment второго snapshot также получает
собственный vector. Это reusable pattern: анализировать независимый снимок,
не воздействуя на исходную inventory.

\begin{shellcode}
cd examples/06_constructors_destructors_and_copy/02_inventory_snapshot
cmake -S . -B build -DCMAKE_BUILD_TYPE=Debug
cmake --build build
./build/inventory_snapshot --root ../../..
ctest --test-dir build --output-on-failure
\end{shellcode}

PASS: 15 entries скопированы construction и assignment, исходный snapshot
остался неизменным, программа печатает \texttt{INVENTORY\_SNAPSHOT=PASS}, оба
CTest имеют code 0 с учётом ожидаемого rejection неверного root.

\subsection{Сильные стороны и цена контроля}

C++ позволяет user-defined type вести себя как обычное value. Programmer при
необходимости точно задаёт creation, copy и destruction; deterministic
destruction создаёт основу управления ресурсами. Composition из string/vector
часто позволяет compiler сгенерировать правильные operations.

Цена: generated shallow copy raw owner создаёт серьёзное UB; special member
rules требуют понимания; deep copy может быть дорогим; ручные destructor и
copy code увеличивают вероятность ошибки. Поэтому ownership должен быть
явным, а обычный современный class по возможности стремится к Rule of Zero.

\subsection{Проверка понимания}

После блока нужно уметь:
1) объяснить начало lifetime и роль constructor в invariant; 2) прочитать member initializer list и назвать настоящий порядок members; 3) объяснить вызов destructor при выходе local object из scope; 4) различить \texttt{T b = a;} и \texttt{b = a;}; 5) описать generated memberwise copy для string/vector; 6) обнаружить owning raw pointer и объяснить shallow-copy double delete; 7) прочитать deep-copy Rule-of-Three implementation и self-assignment; 8) назвать Rule of Three guideline, а не требованием синтаксиса; 9) обосновать Rule of Zero для реального \texttt{CourseBlock}; 10) интерпретировать trace, warning и настоящую ASan diagnostic; 11) оценить deep copy как $O(n)$, не называя generated copy бесплатной; 12) аккуратно учесть copy elision при наблюдении constructor calls.

\subsection{Источники для сверки}

Object lifetime, constructors, mem-initializers, initialization order,
destructors, implicitly declared special members, copying и copy elision
сверены с C++17 working draft N4659, прежде всего clauses 6, 12 и 15.
Педагогическая последовательность constructor/destructor и ownership
сопоставлена с MIT 6.096 Lectures~6 и~8; старые raw-pointer examples сохранены
только как изолированная диагностика, не как modern style. Composition и
standard-library value semantics сверены с локальным \textit{Modern C++
Tutorial}; target-based sanitizer build --- с принятым проектом Modern CMake.
